\section*{Notation and domains of the constructions}
\phantomsection
\label{sec:ii-notacion}
\addcontentsline{toc}{section}{Notation and domains of the constructions}

The table identifies the main symbols by their domain and the
definition introducing them. Subscripts, typography, and operator
context are part of the notation. References point to the
effective definitions; anticipatory occurrences retain that
same meaning.

\begingroup
\small
\setlength{\tabcolsep}{4pt}
\renewcommand{\arraystretch}{1.12}
\begin{longtable}{@{}>{\raggedright\arraybackslash}p{0.19\linewidth}
                    >{\raggedright\arraybackslash}p{0.49\linewidth}
                    >{\raggedright\arraybackslash}p{0.27\linewidth}@{}}
\toprule
Symbol & Object and meaning & Definition and location\\
\midrule
\endfirsthead
\toprule
Symbol & Object and meaning & Definition and location\\
\midrule
\endhead
\bottomrule
\endfoot

\(\mathcal U_t\)
& Transition on the state with memory: selection, transport, and
  record update, in that order.
& \eqref{eq:actualizacion}.\\[3pt]

\(\mathcal R_{12}\), \(R_{\rm sgn}\)
& Twelve-window record with route, orientation, carry, and incidences;
  signed reader with codomain \(\ZZ^{12}\).
& \S\ref{subsec:k-registro-integral};
  composition \eqref{eq:k-lector-firmado}.\\[3pt]

\(A_m,C_m,V_m\)
& Three integer counts per window: residual visits, difference of
  oriented traces, and signed boundary incidence.
& \S\ref{subsec:registro-evaluacion-incidencial},
  definition of the counts.\\[3pt]

\(z_m\)
& Decimal incidence block in \(\{0,\ldots,999\}\):
  \(100[A_m]_{10}+10[C_m]_{10}+[V_m]_{10}\).
& \eqref{eq:registro-bloque-incidencial}.\\[3pt]

\(U\), \(U^{(r)}\)
& Signed integer vector of twelve coordinates and its three
  four-coordinate harmonic blocks.
& \eqref{eq:k-bloques-firmados} and
  \eqref{eq:k-u-firmado}.\\[3pt]

\(K\), \(K_j\)
& Integral record \(K=\mathcal H_{12}U/4\), with fixed phase and
  orientation; coordinate prolongation by \(K_{j+12}=K_j\).
& \S\ref{subsec:k-hadamard}, definition of the integral transformation;
  \S\ref{subsec:alpha-dominio}.\\[3pt]

\(\mathcal K_{\rm ph}\)
& Transfer kernel on \(\mathcal U_6\times C_9\),
  also called \(K_{\rm ph}\) when describing the reduced-memory
  representation.
& \S\ref{ii:tpk-arquitectura}, ``Reduced transfer'';
  return \eqref{ii:tpk-renovacion}.\\[3pt]

\(A_j\), \(c_j\)
& Base-\(1000\) block of \(\alpha\) and integer quotient of
  Euclidean division; index \(j\) runs over the positional expansion.
& \eqref{eq:alpha-recurrencia}.\\[3pt]

\(\mathbf z_m\), \(y_m\)
& Accumulated vector memory in \(\RR^{12}\) and coordinate in the
  moving frame: \(y_m=S^{-m}\mathbf z_m\).
& \S\ref{sec:eta}, increment proposition;
  update \eqref{eq:eta}.\\[3pt]

\(S\), \(R=S^{-1}\)
& Cyclic permutation of the memory basis,
  \(S\mathbf e_m=\mathbf e_{m+1}\); \(R\) denotes its inverse
  within the resolvent development.
& \S\ref{sec:eta};
  \S\ref{subsec:ii-memoria-resolvente}, opening paragraph.\\[3pt]

\(\eta_m\)
& Vector increment produced by the oriented event:
  \(\eta_m=a_mS^{-1}\mathbf e_0\).
& \eqref{eq:events} and \eqref{eq:eta}.\\[3pt]

\(U:X\to X\), \(A\)
& In memory reduction, \(U\) is the state update;
  \(A:V_0\to V_0\) is the retained block of its linear extension
  \(T\delta_x=\delta_{U(x)}\).
& \S\ref{subsec:memoria-resolvente}, ``Memory elimination'';
  Proposition~\ref{mem:prop-schur}.\\[3pt]

\(U:\mathcal H\to\mathcal H\)
& Isometric realization of one-turn transport,
  \(U^*U=I\), in observable compression.
& \S\ref{sec:extension}, ``Observable compression'';
  \eqref{eq:conservacion-ortogonal}.\\[3pt]

\(R=\ZZ/9\ZZ\)
& APP residue ring, with maximal ideal
  \(\mathfrak m=(3)\).
& \S\ref{sec:nucleo}, ``Ternary reduction and multiplicative radical''.\\[3pt]

Local \(R\), \(\eta\)
& In the central block \(B_c=7I+2R\), \(R\) exchanges the two
  coordinates and \(R^2=I\). The rapidity of its normalization is
  \(\eta=\tfrac12\log(9/5)\).
& \S\ref{sec:extension}, ``Local development'':
  \(B_c/\sqrt{45}=\exp(\eta R)\).\\[3pt]

\(\Sigma_H\)
& Dimensionless norm-and-self-scale reader
  \(\Sigma_H=\varphi\alpha^{16}\); its decimal valuation is
  \(\operatorname{ord}_{10}\Sigma_H=-34\).
& \eqref{eq:el-lector-determinantal} and \eqref{eq:el-decada}.\\[3pt]

\pagebreak[4]

\(\mathcal L_S\), \(\mathcal U_S\)
& Oriented action line and positive basis in which
  its sections are expressed.
& \S\ref{sec:action}, ``Action sections and dimensionless
  return quotient''.\\[3pt]

\(H_5\), \(\eta_{\rm ret}\)
& Dimensionless action coefficients before and after return;
  \(\eta_{\rm ret}=H_5-(90/\pi)\mathcal C_\pi(\alpha)\).
& \eqref{eq:el-h5}--\eqref{eq:el-eta}.\\[3pt]

\(\hbar_{\rm pre}\), \(\hbar_{\rm ret}\), \(h_a\)
& Sections \(H_5\,10^{-34}\mathcal U_S\) and
  \(\eta_{\rm ret}\,10^{-34}\mathcal U_S\).
  For \(a\in\{\mathrm{pre},\mathrm{ret}\}\),
  \(h_a=2\pi\hbar_a\) expresses action per turn.
& \eqref{eq:el-secciones-accion};
  \S\ref{sec:action}, passage to action per turn.\\[3pt]

\(R_{\rm act}\), \(\delta_{\rm ret}\)
& Multiplicative ratio \(H_5/\eta_{\rm ret}\) and additive
  difference \(H_5-\eta_{\rm ret}\) of the return coordinates.
& \eqref{eq:el-ratio}.\\[3pt]

\(A_{\deg}\), \(C^*_{\deg}\)
& Angular coordinates \(1000\alpha\) and
  \(2(\eta_{\rm ret}+\alpha)\). Their radian coordinates are
  \(x=\pi A_{\deg}/180\), \(y=\pi C^*_{\deg}/180\).
& \eqref{eq:ii-par-angular};
  \S\ref{sec:ii-angulos}, definition of the charts.\\[3pt]

\(\gamma_{\HMT}\), \(\gamma\)
& Conjugate evaluation of the oriented chain \(c_{12}^{+}\) at
  \(q_\pm\), with angular component \(\pi AC^*/180\), for
  \(0<q_\pm<1\). The positive area increment uses
  \(|\gamma_{\HMT}|\).
& \S\ref{sec:ii-definicion-barbero};
  evaluation \eqref{eq:ii-definicion-barbero-evaluacion}.\\[3pt]

\(\eta_{\rm el}\), \(a_S,b_S\)
& Shape parameter \(\operatorname{artanh}(C^*/A)\) and
  balanced canonical semiaxes of the ellipse, with
  \(a_Sb_S=2\hbar\).
& \S\ref{sec:ellipse}, initial definition.\\[3pt]

\(R_{\rm el}\), \(\tau_{\rm el}\)
& Dimensionless modular radius \(\sqrt{b_S/a_S}\) and modulus
  \(\tau_{\rm el}=iR_{\rm el}^2\).
& \S\ref{sec:ellipse}, reciprocity of the ellipse.\\[3pt]

\(K_a^i\)
& Extrinsic curvature in the normal convention of the geometric
  realization; enters
  \(\mathcal A_a^i=\Gamma_a^i(E)+\gamma K_a^i\).
& \S\ref{sec:ii-costura-barbero}, definition of the connection.\\[3pt]

\(L_*\), \(L\), \(\ell_P\)
& Positive longitudinal reference and length
  \(L=(5759/23040)\alpha_{\HMT}^{16}L_*\).
  Realization R1--R8 determines \(\ell_P^2=L^2\);
  the SI evaluation declares \(L_*=1\,\mathrm m\).
& \eqref{g26:eq:regla-longitudinal-en} and
  \eqref{g26:eq:g-rigido-en}.\\[3pt]

\(R_X\), \(\Lambda_X\), \(E_L\)
& Positive radial reading, operator-valued conjugate length, and
  energy scale \(E_L=\hbar_{\rm ret}c_{\rm int}/L\).
  They satisfy \(R_X\Lambda_X=L^2I\).
& \eqref{g26:eq:lecturas-conjugadas-en} and
  \eqref{g26:eq:productos-radiales-en}.\\[3pt]

\(R_{108}\), horizon \(R\)
& Circular radius \(c_{\rm int}/\omega_{108}\) of the dodecaphasic
  period, equal to \(L\) for the conjugate clock;
  \(R>0\) is the radial variable of the horizon
  realization.
& \S\ref{sec:gravity}, initial definition;
  \eqref{eq:ii-horizonte-normalizado}.\\[3pt]

\(G_a\)
& Gravitational section reciprocal to an action \(\hbar_a>0\):
  \(G_a=c_{\rm int}^{3}L^{2}/\hbar_a\), with \(R_{108}=L\),
  \(c_{\rm int}>0\) and radius coincidence in the convention
  \(r_g=Gm/c_{\rm int}^{2}\).
& \S\ref{sec:ii-definicion-gravedad};
  \eqref{eq:ii-definicion-gravedad-seccion} and
  \eqref{eq:ii-definicion-gravedad-area}.\\[3pt]

\(\mathbb G_{5,a}\)
& Operator \(G_aP^{\mathrm{ico}}_5\) on
  \(V_{12}=\RR[A_5/C_5]\). The orthogonal projector selects
  the summand \(V_5\); section \(G_a\) is common to its five
  components.
& \eqref{ii:pent:eq:gravity-G5};
  Theorem~\ref{ii:pent:thm:gravity-p5-coupling}.\\[3pt]

\(k_B\), \(\Theta_{{\rm clk},a}\)
& Positive linear transducer
  \(k_B:\mathcal L_\Theta\to\mathcal L_E\) and positive thermal
  section, related by
  \(k_B\Theta_{{\rm clk},a}\log3=h_a/(108t_0)\).
  The thermal section and declared bases determine its
  individual coordinate.
& \S\ref{sec:ii-definicion-boltzmann};
  \eqref{eq:ii-kb-aut-par-termico} and
  \eqref{eq:ii-kb-aut-coordenada}.\\

\end{longtable}
\endgroup

In the incidence record, windows carry indices \(1,\ldots,12\).
In the memory update they are numbered from \(0\) to \(11\), and
their prolongation uses \(m\geq0\). The index shift preserves
the same partition into nonadic windows. In angular formulas,
\(A,C^*\) retain the chart declared in each expression; the quotient
\(C^*/A\) is invariant under the common conversion between degrees and radians.
